\subsubsection{Đánh giá kiểm định các cặp mô hình}

\tablename~\ref{tab:regression-stat-tests-mse}
và~\ref{tab:regression-stat-tests-mae} cho kết quả kiểm định các cặp mô hình
tương ứng theo từng chỉ số MSE, MAE.

\import{\tablespath}{regression-stat-tests-mse}
\import{\tablespath}{regression-stat-tests-mae}

Đối với chỉ số MSE, ta thấy mô hình OLS đều cho thấy sự khác biệt với các mô
hình khác với mức ý nghĩa $\alpha = 0.05$. Khi so với kết quả của
\tablename~\ref{tab:regression-models-cv-part1} và
\tablename~\ref{tab:regression-models-cv-part2}, ta biết rằng OLS tốt hơn các
mô hình. Đặc biệt, mặc dù mô hình Polynomial có MSE kém hơn OLS nhưng chưa đủ ý
nghĩa thống kê để chứng minh có sự sai khác với OLS. Điều này có thể do việc
chỉ dùng 10 fold cho cross-validation làm thống kê chưa phân biệt rõ ràng hơn.

Đối với chỉ số MAE, kết quả cho thấy tất cả các cặp mô hình đều có sự sai khác
nhau với mức ý nghĩa $\alpha = 0.05$. Có sự sai khác giữa OLS (dựa trên Normal
Equation) và MBGD do phương pháp tìm nghiệm khác nhau. OLS cho sự ổn định hơn
khi giải trực tiếp nghiệm. Điều này tạo ra sự khác biệt có ý nghĩa thống kê
trên MAE, mặc dù trên MSE không có sự khác biệt.

Kết quả thực nghiệm cho thấy các mô hình WLS, Polynomial và Fourier đạt MAE
thấp hơn so với OLS trong một số thiết lập cross-validation. Điều này cho thấy
việc mở rộng không gian biểu diễn (Polynomial, Fourier) hoặc điều chỉnh trọng
số quan sát (WLS) có thể giúp mô hình khai thác tốt hơn cấu trúc dữ liệu, từ đó
cải thiện độ chính xác dự đoán.

Tuy nhiên, sự cải thiện này không hoàn toàn đồng nhất trên tất cả các phép kiểm
định. Một số cặp so sánh vẫn cho thấy khác biệt có ý nghĩa thống kê, trong khi
một số khác có mức chênh lệch nhỏ, cho thấy hiệu quả cải thiện phụ thuộc vào
phân phối dữ liệu và cấu trúc nhiễu.

Nhìn chung, khi so sánh với các mô hình còn lại, trên các chỉ số còn lại, OLS
thể hiện kết quả ổn định và có tính nhất quán cao trên các phép đo đánh giá.
Mặc dù một số mô hình mở rộng như WLS, Polynomial hay Fourier có thể đạt MAE
tốt hơn trong một số trường hợp cụ thể, sự cải thiện này không mang tính đồng
đều và thường đi kèm với độ biến thiên cao hơn giữa các fold.

Tổng thể, xét trên cả độ chính xác trung bình, tính ổn định và mức độ đơn giản
của mô hình, OLS vẫn cho thấy hiệu quả tốt nhất trong bài toán này, đồng thời
là lựa chọn cân bằng giữa hiệu năng và tính tổng quát hóa.
